Ein Sauerstoffmolekül besteht aus zwei Atomen mit der Atommasse $A_\text{O} = \AOxO$ (molare Masse in g/mol, FoTaBe S.~228) und $u = \uMaO$ (FoTaBe S.~218):
\begin{align}
 m_{\text{O}_2} &= \mOxF = 2\cdot\AOx\cdot\uMa = \mOx
\end{align}
Nach der Absorption bewegen sich Elektron und Molekül gemeinsam. Der Impuls des Elektrons geht auf das Molekül (mit dem Elektron) über:
\begin{align}
 \Delta v &= \dvOF \\
   &= \frac{\pEl}{\mOx+\mel} \\
   &= \dvO \approx \result{\dvOP{3}{2}}
\end{align}
Die Elektronenmasse \mel\ ist gegenüber der Molekülmasse vernachlässigbar.
